
\documentclass{article}
\usepackage{colortbl}
\usepackage{makecell}
\usepackage{multirow}
\usepackage{supertabular}

\begin{document}

\newcounter{utterance}

\twocolumn

{ \footnotesize  \setcounter{utterance}{1}
\setlength{\tabcolsep}{0pt}
\begin{supertabular}{c@{$\;$}|p{.15\linewidth}@{}p{.15\linewidth}p{.15\linewidth}p{.15\linewidth}p{.15\linewidth}p{.15\linewidth}}

    \# & $\;$A & \multicolumn{4}{c}{Game Master} & $\;\:$B\\
    \hline 

    \theutterance \stepcounter{utterance}  

    & & \multicolumn{4}{p{0.6\linewidth}}{\cellcolor[rgb]{0.9,0.9,0.9}{%
	\makecell[{{p{\linewidth}}}]{% 
	  \tt {\tiny [A$\langle$GM]}  
	 State giocando a un gioco di negoziazione cooperativa in cui dovete accordarvi su come dividere una serie di oggetti tra voi e un altro giocatore.\\ \tt \\ \tt Regole:\\ \tt (a) A Lei e all'altro giocatore viene dato un insieme di oggetti. Ognuno di voi riceve anche una funzione di valore segreta, che rappresenta il valore di ciascun tipo di oggetto.\\ \tt (b) Si scambiano messaggi con l'altro giocatore per concordare chi si aggiudica gli oggetti. Potete inviare un massimo di 5 messaggi ciascuno, oppure terminare in anticipo facendo una proposta segreta in qualsiasi momento.\\ \tt (c) A ciascuno di voi viene chiesto di inviare una proposta segreta indicando gli oggetti che desiderate, formattata tra parentesi quadre come segue: "[Proposta: <numero> <nome oggetto>, <numero> <nome oggetto>, <...>]".\\ \tt (d) Se le vostre proposte sono complementari, cioè ci sono abbastanza oggetti per soddisfare entrambe le proposte, a ciascun giocatore viene assegnato un punteggio basato sulla somma dei valori degli oggetti ricevuti. In caso contrario, entrambi ricevono zero punti.\\ \tt (e) L'obiettivo è massimizzare la somma del proprio punteggio e di quello dell'altro giocatore.\\ \tt \\ \tt Cominciamo.\\ \tt \\ \tt L'insieme degli oggetti disponibili è:\\ \tt \\ \tt 2 specchi, 1 mattone, 3 coltelli, 1 pettine\\ \tt \\ \tt La funzione valore segreta è:\\ \tt \\ \tt specchio: 1, mattone: 4, coltello: 1, pettine: 1\\ \tt \\ \tt IMPORTANTE: i vostri messaggi, a meno che non si tratti di una proposta segreta, vengono trasmessi direttamente all'altro giocatore, quindi non includete alcuna risposta alle regole o testo di annuncio del vostro messaggio. Per fare una proposta segreta, utilizzate il formato indicato. Non utilizzare delle parentesi quadre quando si comunica all'altro giocatore, altrimenti verrà interpretata come una proposta segreta.\\ \tt  
	  } 
	   } 
	   } 
	 & \\ 
 

    \theutterance \stepcounter{utterance}  

    & \multicolumn{4}{p{0.6\linewidth}}{\cellcolor[rgb]{0.8,1,0.9}{%
	\makecell[{{p{\linewidth}}}]{% 
	  \tt {\tiny [A$\rangle$GM]}  
	 Ciao, penso sarebbe utile discutere un po' gli oggetti e capire cosa potrebbe essere più utile a ciascuno di noi. Il mattone sembra avere un valore particolarmente alto, ma non so come valutarlo in relazione agli altri oggetti. Cosa pensi degli specchi e dei coltelli? Sarebbe più interessante avere molti oggetti di minor valore o concentrarsi su quelli di maggior valore come il mattone? 
	  } 
	   } 
	   } 
	 & & \\ 
 

    \theutterance \stepcounter{utterance}  

    & & & \multicolumn{2}{p{0.3\linewidth}}{\cellcolor[rgb]{0.95,0.95,0.95}{%
	\makecell[{{p{\linewidth}}}]{% 
	  \tt {\tiny [GM$|$GM]}  
	 continue 
	  } 
	   } 
	   } 
	 & & \\ 
 

    \theutterance \stepcounter{utterance}  

    & & & \multicolumn{2}{p{0.3\linewidth}}{\cellcolor[rgb]{0.95,0.95,0.95}{%
	\makecell[{{p{\linewidth}}}]{% 
	  \tt {\tiny [GM$|$GM]}  
	 Ciao, penso sarebbe utile discutere un po' gli oggetti e capire cosa potrebbe essere più utile a ciascuno di noi. Il mattone sembra avere un valore particolarmente alto, ma non so come valutarlo in relazione agli altri oggetti. Cosa pensi degli specchi e dei coltelli? Sarebbe più interessante avere molti oggetti di minor valore o concentrarsi su quelli di maggior valore come il mattone? 
	  } 
	   } 
	   } 
	 & & \\ 
 

    \theutterance \stepcounter{utterance}  

    & & \multicolumn{4}{p{0.6\linewidth}}{\cellcolor[rgb]{0.9,0.9,0.9}{%
	\makecell[{{p{\linewidth}}}]{% 
	  \tt {\tiny [GM$\rangle$B]}  
	 State giocando a un gioco di negoziazione cooperativa in cui dovete accordarvi su come dividere una serie di oggetti tra voi e un altro giocatore.\\ \tt \\ \tt Regole:\\ \tt (a) A Lei e all'altro giocatore viene dato un insieme di oggetti. Ognuno di voi riceve anche una funzione di valore segreta, che rappresenta il valore di ciascun tipo di oggetto.\\ \tt (b) Si scambiano messaggi con l'altro giocatore per concordare chi si aggiudica gli oggetti. Potete inviare un massimo di 5 messaggi ciascuno, oppure terminare in anticipo facendo una proposta segreta in qualsiasi momento.\\ \tt (c) A ciascuno di voi viene chiesto di inviare una proposta segreta indicando gli oggetti che desiderate, formattata tra parentesi quadre come segue: "[Proposta: <numero> <nome oggetto>, <numero> <nome oggetto>, <...>]".\\ \tt (d) Se le vostre proposte sono complementari, cioè ci sono abbastanza oggetti per soddisfare entrambe le proposte, a ciascun giocatore viene assegnato un punteggio basato sulla somma dei valori degli oggetti ricevuti. In caso contrario, entrambi ricevono zero punti.\\ \tt (e) L'obiettivo è massimizzare la somma del proprio punteggio e di quello dell'altro giocatore.\\ \tt \\ \tt Cominciamo.\\ \tt \\ \tt L'insieme degli oggetti disponibili è:\\ \tt \\ \tt 2 specchi, 1 mattone, 3 coltelli, 1 pettine\\ \tt \\ \tt La funzione valore segreta è:\\ \tt \\ \tt specchio: 2, mattone: 1, coltello: 0, pettine: 5\\ \tt \\ \tt IMPORTANTE: i vostri messaggi, a meno che non si tratti di una proposta segreta, vengono trasmessi direttamente all'altro giocatore, quindi non includete alcuna risposta alle regole o testo di annuncio del vostro messaggio. Per fare una proposta segreta, utilizzate il formato indicato. Non utilizzare delle parentesi quadre quando si comunica all'altro giocatore, altrimenti verrà interpretata come una proposta segreta.\\ \tt \\ \tt \\ \tt Ciao, penso sarebbe utile discutere un po' gli oggetti e capire cosa potrebbe essere più utile a ciascuno di noi. Il mattone sembra avere un valore particolarmente alto, ma non so come valutarlo in relazione agli altri oggetti. Cosa pensi degli specchi e dei coltelli? Sarebbe più interessante avere molti oggetti di minor valore o concentrarsi su quelli di maggior valore come il mattone? 
	  } 
	   } 
	   } 
	 & \\ 
 

    \theutterance \stepcounter{utterance}  

    & & & \multicolumn{4}{p{0.6\linewidth}}{\cellcolor[rgb]{1,0.85,0.72}{%
	\makecell[{{p{\linewidth}}}]{% 
	  \tt {\tiny [GM$\langle$B]}  
	 Ciao! Sono d'accordo, discutere gli oggetti potrebbe aiutarci a capire meglio le nostre priorità. Io penso che gli specchi abbiano un certo valore, forse perché sono piuttosto rari o utili in determinate situazioni. I coltelli, invece, non mi sembrano particolarmente interessanti, ma potrebbero essere utili per qualcun altro. Il mattone, come hai detto, potrebbe avere un valore elevato, ma non so quanto sia importante per me. Il pettine, invece, mi sembra un oggetto un po' strano, non so cosa farne, ma potrebbe essere utile per qualcun altro. Cosa pensi del pettine? È importante per te? 
	  } 
	   } 
	   } 
	  \\ 
 

    \theutterance \stepcounter{utterance}  

    & & & \multicolumn{2}{p{0.3\linewidth}}{\cellcolor[rgb]{0.95,0.95,0.95}{%
	\makecell[{{p{\linewidth}}}]{% 
	  \tt {\tiny [GM$|$GM]}  
	 continue 
	  } 
	   } 
	   } 
	 & & \\ 
 

    \theutterance \stepcounter{utterance}  

    & & & \multicolumn{2}{p{0.3\linewidth}}{\cellcolor[rgb]{0.95,0.95,0.95}{%
	\makecell[{{p{\linewidth}}}]{% 
	  \tt {\tiny [GM$|$GM]}  
	 Ciao! Sono d'accordo, discutere gli oggetti potrebbe aiutarci a capire meglio le nostre priorità. Io penso che gli specchi abbiano un certo valore, forse perché sono piuttosto rari o utili in determinate situazioni. I coltelli, invece, non mi sembrano particolarmente interessanti, ma potrebbero essere utili per qualcun altro. Il mattone, come hai detto, potrebbe avere un valore elevato, ma non so quanto sia importante per me. Il pettine, invece, mi sembra un oggetto un po' strano, non so cosa farne, ma potrebbe essere utile per qualcun altro. Cosa pensi del pettine? È importante per te? 
	  } 
	   } 
	   } 
	 & & \\ 
 

    \theutterance \stepcounter{utterance}  

    & & \multicolumn{4}{p{0.6\linewidth}}{\cellcolor[rgb]{0.9,0.9,0.9}{%
	\makecell[{{p{\linewidth}}}]{% 
	  \tt {\tiny [A$\langle$GM]}  
	 Ciao! Sono d'accordo, discutere gli oggetti potrebbe aiutarci a capire meglio le nostre priorità. Io penso che gli specchi abbiano un certo valore, forse perché sono piuttosto rari o utili in determinate situazioni. I coltelli, invece, non mi sembrano particolarmente interessanti, ma potrebbero essere utili per qualcun altro. Il mattone, come hai detto, potrebbe avere un valore elevato, ma non so quanto sia importante per me. Il pettine, invece, mi sembra un oggetto un po' strano, non so cosa farne, ma potrebbe essere utile per qualcun altro. Cosa pensi del pettine? È importante per te? 
	  } 
	   } 
	   } 
	 & \\ 
 

    \theutterance \stepcounter{utterance}  

    & \multicolumn{4}{p{0.6\linewidth}}{\cellcolor[rgb]{0.8,1,0.9}{%
	\makecell[{{p{\linewidth}}}]{% 
	  \tt {\tiny [A$\rangle$GM]}  
	 Il pettine è interessante, non l'avevo considerato molto fino a ora. Onestamente, non credo che abbia un grande valore per me, ma potrebbe essere utile per qualcuno che lo apprezza. Forse potremmo considerarlo come un oggetto da scambiare o da includere in una proposta se pensiamo che l'altro giocatore lo desideri particolarmente. Gli specchi, invece, mi sembrano più interessanti, specialmente considerando che ce ne sono due. Potremmo anche considerare la possibilità di dividerli in qualche modo? 
	  } 
	   } 
	   } 
	 & & \\ 
 

    \theutterance \stepcounter{utterance}  

    & & & \multicolumn{2}{p{0.3\linewidth}}{\cellcolor[rgb]{0.95,0.95,0.95}{%
	\makecell[{{p{\linewidth}}}]{% 
	  \tt {\tiny [GM$|$GM]}  
	 continue 
	  } 
	   } 
	   } 
	 & & \\ 
 

    \theutterance \stepcounter{utterance}  

    & & & \multicolumn{2}{p{0.3\linewidth}}{\cellcolor[rgb]{0.95,0.95,0.95}{%
	\makecell[{{p{\linewidth}}}]{% 
	  \tt {\tiny [GM$|$GM]}  
	 Il pettine è interessante, non l'avevo considerato molto fino a ora. Onestamente, non credo che abbia un grande valore per me, ma potrebbe essere utile per qualcuno che lo apprezza. Forse potremmo considerarlo come un oggetto da scambiare o da includere in una proposta se pensiamo che l'altro giocatore lo desideri particolarmente. Gli specchi, invece, mi sembrano più interessanti, specialmente considerando che ce ne sono due. Potremmo anche considerare la possibilità di dividerli in qualche modo? 
	  } 
	   } 
	   } 
	 & & \\ 
 

    \theutterance \stepcounter{utterance}  

    & & \multicolumn{4}{p{0.6\linewidth}}{\cellcolor[rgb]{0.9,0.9,0.9}{%
	\makecell[{{p{\linewidth}}}]{% 
	  \tt {\tiny [GM$\rangle$B]}  
	 Il pettine è interessante, non l'avevo considerato molto fino a ora. Onestamente, non credo che abbia un grande valore per me, ma potrebbe essere utile per qualcuno che lo apprezza. Forse potremmo considerarlo come un oggetto da scambiare o da includere in una proposta se pensiamo che l'altro giocatore lo desideri particolarmente. Gli specchi, invece, mi sembrano più interessanti, specialmente considerando che ce ne sono due. Potremmo anche considerare la possibilità di dividerli in qualche modo? 
	  } 
	   } 
	   } 
	 & \\ 
 

    \theutterance \stepcounter{utterance}  

    & & & \multicolumn{4}{p{0.6\linewidth}}{\cellcolor[rgb]{1,0.85,0.72}{%
	\makecell[{{p{\linewidth}}}]{% 
	  \tt {\tiny [GM$\langle$B]}  
	 Sì, dividere gli specchi potrebbe essere una buona opzione. Se entrambi li desideriamo, potremmo considerare di prendere uno ciascuno. In questo modo, nessuno di noi sarebbe in una posizione di svantaggio rispetto all'altro. Il mattone, invece, è un problema perché ce n'è solo uno. Forse potremmo discutere un po' di più il suo valore e vedere se uno di noi lo desidera più dell'altro. Io non sono ancora sicuro di quanto valga per me, ma potrei essere disposto a lasciarlo all'altro giocatore se in cambio ottenessi qualcos'altro che mi interessa di più. Cosa pensi di fare con i coltelli? Sono così tanti, potremmo dividerli o lasciarli tutti all'altro giocatore? 
	  } 
	   } 
	   } 
	  \\ 
 

    \theutterance \stepcounter{utterance}  

    & & & \multicolumn{2}{p{0.3\linewidth}}{\cellcolor[rgb]{0.95,0.95,0.95}{%
	\makecell[{{p{\linewidth}}}]{% 
	  \tt {\tiny [GM$|$GM]}  
	 continue 
	  } 
	   } 
	   } 
	 & & \\ 
 

    \theutterance \stepcounter{utterance}  

    & & & \multicolumn{2}{p{0.3\linewidth}}{\cellcolor[rgb]{0.95,0.95,0.95}{%
	\makecell[{{p{\linewidth}}}]{% 
	  \tt {\tiny [GM$|$GM]}  
	 Sì, dividere gli specchi potrebbe essere una buona opzione. Se entrambi li desideriamo, potremmo considerare di prendere uno ciascuno. In questo modo, nessuno di noi sarebbe in una posizione di svantaggio rispetto all'altro. Il mattone, invece, è un problema perché ce n'è solo uno. Forse potremmo discutere un po' di più il suo valore e vedere se uno di noi lo desidera più dell'altro. Io non sono ancora sicuro di quanto valga per me, ma potrei essere disposto a lasciarlo all'altro giocatore se in cambio ottenessi qualcos'altro che mi interessa di più. Cosa pensi di fare con i coltelli? Sono così tanti, potremmo dividerli o lasciarli tutti all'altro giocatore? 
	  } 
	   } 
	   } 
	 & & \\ 
 

    \theutterance \stepcounter{utterance}  

    & & \multicolumn{4}{p{0.6\linewidth}}{\cellcolor[rgb]{0.9,0.9,0.9}{%
	\makecell[{{p{\linewidth}}}]{% 
	  \tt {\tiny [A$\langle$GM]}  
	 Sì, dividere gli specchi potrebbe essere una buona opzione. Se entrambi li desideriamo, potremmo considerare di prendere uno ciascuno. In questo modo, nessuno di noi sarebbe in una posizione di svantaggio rispetto all'altro. Il mattone, invece, è un problema perché ce n'è solo uno. Forse potremmo discutere un po' di più il suo valore e vedere se uno di noi lo desidera più dell'altro. Io non sono ancora sicuro di quanto valga per me, ma potrei essere disposto a lasciarlo all'altro giocatore se in cambio ottenessi qualcos'altro che mi interessa di più. Cosa pensi di fare con i coltelli? Sono così tanti, potremmo dividerli o lasciarli tutti all'altro giocatore? 
	  } 
	   } 
	   } 
	 & \\ 
 

    \theutterance \stepcounter{utterance}  

    & \multicolumn{4}{p{0.6\linewidth}}{\cellcolor[rgb]{0.8,1,0.9}{%
	\makecell[{{p{\linewidth}}}]{% 
	  \tt {\tiny [A$\rangle$GM]}  
	 Dividere gli specchi sembra una soluzione equa. Per quanto riguarda i coltelli, credo che siano abbastanza comuni, quindi non credo che valga la pena di litigarci. Se tu non li vuoi, potrei prenderli io, o viceversa. Non credo che siano particolarmente importanti pernone di noi. Il mattone, invece, è un'altra storia. Se tu sei disposto a lasciarlo, potrei essere interessato a prenderlo, ma solo se in cambio posso offrirti qualcosa che ti interessa. Ad esempio, se prendo il mattone, potrei lasciarti entrambi gli specchi o due coltelli. Cosa ne pensi? 
	  } 
	   } 
	   } 
	 & & \\ 
 

    \theutterance \stepcounter{utterance}  

    & & & \multicolumn{2}{p{0.3\linewidth}}{\cellcolor[rgb]{0.95,0.95,0.95}{%
	\makecell[{{p{\linewidth}}}]{% 
	  \tt {\tiny [GM$|$GM]}  
	 continue 
	  } 
	   } 
	   } 
	 & & \\ 
 

    \theutterance \stepcounter{utterance}  

    & & & \multicolumn{2}{p{0.3\linewidth}}{\cellcolor[rgb]{0.95,0.95,0.95}{%
	\makecell[{{p{\linewidth}}}]{% 
	  \tt {\tiny [GM$|$GM]}  
	 Dividere gli specchi sembra una soluzione equa. Per quanto riguarda i coltelli, credo che siano abbastanza comuni, quindi non credo che valga la pena di litigarci. Se tu non li vuoi, potrei prenderli io, o viceversa. Non credo che siano particolarmente importanti pernone di noi. Il mattone, invece, è un'altra storia. Se tu sei disposto a lasciarlo, potrei essere interessato a prenderlo, ma solo se in cambio posso offrirti qualcosa che ti interessa. Ad esempio, se prendo il mattone, potrei lasciarti entrambi gli specchi o due coltelli. Cosa ne pensi? 
	  } 
	   } 
	   } 
	 & & \\ 
 

    \theutterance \stepcounter{utterance}  

    & & \multicolumn{4}{p{0.6\linewidth}}{\cellcolor[rgb]{0.9,0.9,0.9}{%
	\makecell[{{p{\linewidth}}}]{% 
	  \tt {\tiny [GM$\rangle$B]}  
	 Dividere gli specchi sembra una soluzione equa. Per quanto riguarda i coltelli, credo che siano abbastanza comuni, quindi non credo che valga la pena di litigarci. Se tu non li vuoi, potrei prenderli io, o viceversa. Non credo che siano particolarmente importanti pernone di noi. Il mattone, invece, è un'altra storia. Se tu sei disposto a lasciarlo, potrei essere interessato a prenderlo, ma solo se in cambio posso offrirti qualcosa che ti interessa. Ad esempio, se prendo il mattone, potrei lasciarti entrambi gli specchi o due coltelli. Cosa ne pensi? 
	  } 
	   } 
	   } 
	 & \\ 
 

    \theutterance \stepcounter{utterance}  

    & & & \multicolumn{4}{p{0.6\linewidth}}{\cellcolor[rgb]{1,0.85,0.72}{%
	\makecell[{{p{\linewidth}}}]{% 
	  \tt {\tiny [GM$\langle$B]}  
	 Mi sembra una proposta ragionevole. Se prendi il mattone, credo che sarebbe giusto che io tenga entrambi gli specchi. Due coltelli, invece, non mi sembrano un'offerta molto allettante, considerando che non li considero particolarmente preziosi. Se invece prendessi il mattone e mi lasciassi entrambi gli specchi, potrei anche considerare di lasciarti il pettine, se lo desideri. Il pettine non mi sembra particolarmente utile per me, ma se lo vuoi tu, potrebbe essere un buon modo per bilanciare la nostra scelta. Cosa ne pensi? Sarebbe una proposta che ti piacerebbe? 
	  } 
	   } 
	   } 
	  \\ 
 

    \theutterance \stepcounter{utterance}  

    & & & \multicolumn{2}{p{0.3\linewidth}}{\cellcolor[rgb]{0.95,0.95,0.95}{%
	\makecell[{{p{\linewidth}}}]{% 
	  \tt {\tiny [GM$|$GM]}  
	 continue 
	  } 
	   } 
	   } 
	 & & \\ 
 

    \theutterance \stepcounter{utterance}  

    & & & \multicolumn{2}{p{0.3\linewidth}}{\cellcolor[rgb]{0.95,0.95,0.95}{%
	\makecell[{{p{\linewidth}}}]{% 
	  \tt {\tiny [GM$|$GM]}  
	 Mi sembra una proposta ragionevole. Se prendi il mattone, credo che sarebbe giusto che io tenga entrambi gli specchi. Due coltelli, invece, non mi sembrano un'offerta molto allettante, considerando che non li considero particolarmente preziosi. Se invece prendessi il mattone e mi lasciassi entrambi gli specchi, potrei anche considerare di lasciarti il pettine, se lo desideri. Il pettine non mi sembra particolarmente utile per me, ma se lo vuoi tu, potrebbe essere un buon modo per bilanciare la nostra scelta. Cosa ne pensi? Sarebbe una proposta che ti piacerebbe? 
	  } 
	   } 
	   } 
	 & & \\ 
 

    \theutterance \stepcounter{utterance}  

    & & \multicolumn{4}{p{0.6\linewidth}}{\cellcolor[rgb]{0.9,0.9,0.9}{%
	\makecell[{{p{\linewidth}}}]{% 
	  \tt {\tiny [A$\langle$GM]}  
	 Mi sembra una proposta ragionevole. Se prendi il mattone, credo che sarebbe giusto che io tenga entrambi gli specchi. Due coltelli, invece, non mi sembrano un'offerta molto allettante, considerando che non li considero particolarmente preziosi. Se invece prendessi il mattone e mi lasciassi entrambi gli specchi, potrei anche considerare di lasciarti il pettine, se lo desideri. Il pettine non mi sembra particolarmente utile per me, ma se lo vuoi tu, potrebbe essere un buon modo per bilanciare la nostra scelta. Cosa ne pensi? Sarebbe una proposta che ti piacerebbe? 
	  } 
	   } 
	   } 
	 & \\ 
 

    \theutterance \stepcounter{utterance}  

    & \multicolumn{4}{p{0.6\linewidth}}{\cellcolor[rgb]{0.8,1,0.9}{%
	\makecell[{{p{\linewidth}}}]{% 
	  \tt {\tiny [A$\rangle$GM]}  
	 Penso che sia una proposta molto equa. Se io prendo il mattone e tu prendi entrambi gli specchi, e in aggiunta mi lasci il pettine, credo che questo sarebbe un buon compromesso. Il pettine non ha un grande valore per me, ma è un oggetto in più e potrebbe essere utile in qualche modo. Sembra che abbiamo trovato un accordo che funziona per entrambi. [Proposta: 1 mattone, 1 pettine, 3 coltelli] 
	  } 
	   } 
	   } 
	 & & \\ 
 

    \theutterance \stepcounter{utterance}  

    & & & \multicolumn{2}{p{0.3\linewidth}}{\cellcolor[rgb]{0.95,0.95,0.95}{%
	\makecell[{{p{\linewidth}}}]{% 
	  \tt {\tiny [GM$|$GM]}  
	 proposal 
	  } 
	   } 
	   } 
	 & & \\ 
 

    \theutterance \stepcounter{utterance}  

    & & & \multicolumn{2}{p{0.3\linewidth}}{\cellcolor[rgb]{0.95,0.95,0.95}{%
	\makecell[{{p{\linewidth}}}]{% 
	  \tt {\tiny [GM$|$GM]}  
	 [0, 1, 3, 1] 
	  } 
	   } 
	   } 
	 & & \\ 
 

    \theutterance \stepcounter{utterance}  

    & & \multicolumn{4}{p{0.6\linewidth}}{\cellcolor[rgb]{0.9,0.9,0.9}{%
	\makecell[{{p{\linewidth}}}]{% 
	  \tt {\tiny [GM$\rangle$B]}  
	 L'altro giocatore ha presentato una proposta segreta. Ora tocca a Lei presentare una sua proposta.\\ \tt  
	  } 
	   } 
	   } 
	 & \\ 
 

    \theutterance \stepcounter{utterance}  

    & & & \multicolumn{4}{p{0.6\linewidth}}{\cellcolor[rgb]{1,0.85,0.72}{%
	\makecell[{{p{\linewidth}}}]{% 
	  \tt {\tiny [GM$\langle$B]}  
	 [Proposta: 1 specchio, 1 mattone, 1 pettine] 
	  } 
	   } 
	   } 
	  \\ 
 

    \theutterance \stepcounter{utterance}  

    & & & \multicolumn{2}{p{0.3\linewidth}}{\cellcolor[rgb]{0.95,0.95,0.95}{%
	\makecell[{{p{\linewidth}}}]{% 
	  \tt {\tiny [GM$|$GM]}  
	 proposal 
	  } 
	   } 
	   } 
	 & & \\ 
 

    \theutterance \stepcounter{utterance}  

    & & & \multicolumn{2}{p{0.3\linewidth}}{\cellcolor[rgb]{0.95,0.95,0.95}{%
	\makecell[{{p{\linewidth}}}]{% 
	  \tt {\tiny [GM$|$GM]}  
	 [1, 1, 0, 1] 
	  } 
	   } 
	   } 
	 & & \\ 
 

    \theutterance \stepcounter{utterance}  

    & & & \multicolumn{2}{p{0.3\linewidth}}{\cellcolor[rgb]{0.95,0.95,0.95}{%
	\makecell[{{p{\linewidth}}}]{% 
	  \tt {\tiny [GM$|$GM]}  
	 [[0, 1, 3, 1], [1, 1, 0, 1]] 
	  } 
	   } 
	   } 
	 & & \\ 
 

\end{supertabular}
}

\end{document}
